\begin{figure}[t]
  \centering
  \fbox{
    \begin{minipage}[c][4cm][c]{0.95\columnwidth}
      \centering
      Visualization examples of the score maps of the initial pose sampling module $\initpose$.
    \end{minipage}
  }
  \caption{Score maps of the initial pose sampling module $\initpose$.}
  \label{fig:score_map}
\end{figure}


\section{Sequence Planning} \label{sec:sequence_planning}

In this section, we introduce the remaining modules $\initpose$ and $\simul$ that are used to generate candidate poses and simulate the stability of the scene, respectively.
Then, we present the Monte Carlo Tree Search (MCTS) algorithm that uses these modules to explore the space of possible stacking sequences and find an optimal sequence of object placements.

\subsection{Initial Pose Sampling}
The initial pose sampler generates a candidate object, orientation, and planar
translation from the current state $\state_k=(\scene_k,\inventory_k,\target)$.
Let $\mathcal{G}_{i,R}(\state_k)$ denote the feasible planar sampling grid for
object $i$ and orientation $R$, and let
$\hat{q}_{k+1}=(R,(\planar,z_{i,R}(\planar;\state_k)))$.
The sampler is defined by
\begin{equation*}
    \begin{aligned}
        \initpose(i,R,\planar\mid\state_k)
        &= \pi_{xy}(\planar\mid i,R,\state_k)
           \pi_\text{unif}(i,R\mid\state_k),\\
        \pi_{xy}(\planar\mid i,R,\state_k)
        &\propto
        \exp\!\left(\score(\planar\mid i,R,\state_k)/\tau\right).
    \end{aligned}
\end{equation*}
where $\pi_\text{unif}(i,R\mid\state_k)$ is uniform over
$i\in\inventory_k$ and the finite rotation set $\mathcal{R}$.
The distribution $\pi_{xy}$ is defined over
$\planar\in\mathcal{G}_{i,R}(\state_k)$, and $\tau>0$ controls the preference
for higher-scoring planar locations. The initial height $z_{i,R}$ is computed
from the scene geometry, while the full pose is subsequently refined by the
pose optimization module $\poseopt$.
For brevity, the dependence of the score on $\state_k$ is omitted below.
The proposal score $\score(\planar\mid i_{k+1},R)$ is formulated as a weighted sum of several geometric features that are computed based on the scene $\scene_k$, the target shape $\target$, and the object $i_{k+1}$:
\begin{equation*}
  \begin{split}
    \score(\planar\mid i_{k+1},R) ={}& w_\text{s}\score_\text{s}(\planar\mid i_{k+1},R)+w_\text{f}\score_\text{f}(\planar\mid i_{k+1},R) \\
    &+w_\text{b}\score_\text{b}(\planar\mid i_{k+1},R)-w_\text{h}\score_\text{h}(\planar\mid i_{k+1},R),
  \end{split}
%   \label{eq:initial_pose_score}
\end{equation*}
where $\score_\text{s}$ is the support area score map, $\score_\text{f}$ is the filled floor area score map, $\score_\text{b}$ is the boundary proximity score map, and $\score_\text{h}$ is the height score map, respectively with weight parameters $w_\text{h}$, $w_\text{s}$, $w_\text{f}, w_\text{b} > 0$.

Intuitively, the score $\score_\text{s}$ encourages the object to be placed on a larger support area, the score $\score_\text{f}$ encourages the object to fill the empty floor area, the score $\score_\text{b}$ encourages to fill the boundary of the target shape $\target$ first, and the last score $\score_\text{h}$ encourages the object to be placed at a lower height to improve stability.
The visualization examples of the score maps are shown in Fig. \ref{fig:score_map} and the details of the score map computation will be described in the appendix.

\begin{algorithm}[t]
  \caption{Initial pose sampling $\initpose(\state_k)$}
  \label{alg:initial_pose_sampling}
  \begin{algorithmic}[1]
    \Require State $\state_k=(\scene_k,\inventory_k,\target)$, rotation set $\mathcal{R}$
    \State Sample $i_{k+1}\sim\operatorname{Unif}(\inventory_k)$ and
           $R\sim\operatorname{Unif}(\mathcal{R})$
    \State Construct the feasible planar grid
           $\mathcal{G}_{i_{k+1},R}(\state_k)$
    \State Evaluate $\score(\planar\mid i_{k+1},R,\state_k)$ for all
           $\planar\in\mathcal{G}_{i_{k+1},R}(\state_k)$
    \State Sample $\planar$ from $\pi_{xy}(\planar\mid i_{k+1},R)$
    \State Set $\hat{q}_{k+1}=(R,(\planar,z_{i_{k+1},R}(\planar;\state_k)))$
    \State \Return $(i_{k+1},\hat{q}_{k+1})$
  \end{algorithmic}
\end{algorithm}

\subsection{Simulation Module}

After $\poseopt$ refines the candidate pose $\hat{q}_{k+1}$, the simulation
module $\simul$ evaluates the deterministic transition
$\state_{k+1}=\transition(\state_k,\action_k)$. It first simulates the new
object with the existing scene fixed to refine the placement, and then
simulates the entire scene to verify global stability.

Both stages use early termination. The simulation stops when all placed
objects have sufficiently small kinetic energy or when accumulated motion
energy indicates instability. For object $\object_i$, we define
\begin{align*}
    E_{\text{kin},i}(t) = \frac{1}{2m_i}\mathcal{V}_i^\top M_i \mathcal{V}_i, \quad 
    E_{\text{acc},i}(t) = \sum_{s=0}^{t} \sqrt{E_{\text{kin},i}(s)} \Delta t
\end{align*}
where $\mathcal{V}_i$ is the spatial velocity, $M_i$ and $m_i$ are the mass
matrix and mass, respectively, and $\Delta t$ is the simulation time step.
The early-termination condition is
\begin{equation}
    \max_{i \in \placed_{k+1}} E_{\text{kin},i}(t) < \epsilon_\text{kin}
    \quad\text{or}\quad
    \max_{i \in \placed_{k+1}} E_{\text{acc},i}(t) > \epsilon_\text{acc}
\end{equation}
with the thresholds $\epsilon_\text{kin}, \epsilon_\text{acc} > 0$.
We implement $\simul$ in CUDA using the GPU-accelerated subADMM algorithm
\cite{ji2025gpu}.

\subsection{Reward Design}

The reward function is designed to encourage the agent to place objects in a way that maximizes the filled volume of the target shape $\target$ while ensuring the stability of the scene and formulated as follows:
\begin{equation*}
    \begin{split}
        r(\state_k, \action_k) 
        ={}& w_\target\frac{\vol\left(\overline{\scene}_{k+1}\right) - \vol\left(\overline{\scene}_k\right)}{\vol(\target)} + w_\object\frac{\vol\left(O_{i_{k+1}}\cap\target\right)}{\vol\left(O_{i_{k+1}}\right)} \\
        &- w_H\left(\max_{\planar}H_{\scene_{k+1}}(\planar) - \mathbb{E}_{\planar}[H_{\scene_{k+1}}(\planar)]\right) \\
        & - w_D\mathbb{E}_{\mathcal{V}\sim\mathcal{N}(0,\Sigma)}\left[\Delta q_{k+1}(\mathcal{V})\right]- w_S f_\text{eq}
    \end{split}
    \label{eq:reward_function}
\end{equation*}
where $w_\target, w_\object, w_H, w_S, w_D > 0$ are the weight parameters for each reward term, $\overline{\scene}_k := T \cap\bigcup_{i\in\placed_k} O_i(q_i)$ is the intersection of the target shape $\target$ and the union of the placed objects in the scene $\scene_k$, and $H_{\scene_{k+1}}(\planar)$ is the uppper surface height of the scene $\scene_{k+1}$ at planar location $\planar = (x,y)$.
The first term and the second term encourage the agent to place objects that fill the target shape $\target$ and maximize the filled volume of the target shape $\target$, respectively.
The third term penalizes the height difference between the maximum height and the average height of the scene $\scene_{k+1}$ to encourage a more uniform distribution of objects in the scene.
The fourth term penalizes the expected pose deviation $\Delta q_{k+1}(\mathcal{V})$ of the newly placed object $\object_{i_{k+1}}$ under random perturbations $\mathcal{V}$ sampled from a Gaussian distribution $\mathcal{N}(0,\Sigma)$ to encourage the agent to place objects in a more stable manner.
The last term penalizes the violation of the static equilibrium constraints $f_\text{eq}$ from $\poseopt$ to encourage the agent to place objects in a physically feasible manner.

\subsection{Monte Carlo Tree Search}
We combine $\initpose$, $\poseopt$, and $\simul$ in a batched,
receding-horizon MCTS planner over the MDP defined in
Sec.~\ref{sec:overview}. Each search considers a finite horizon of future
placements, after which only the selected immediate placement is executed.
Since evaluating a placement requires both pose optimization and physical
simulation, we select multiple search paths before evaluating their candidate
placements in parallel, following the batched planning approach
of~\cite{huang2022pmbs}. Algorithm~\ref{alg:mcts} summarizes the procedure.

At the beginning of each planning step, the tree $\tree$ is initialized or updated to match the current state $\state_k$. 
Root preparation combines fresh proposals from $\initpose$ with mutations of promising candidates and evaluates them through the same optimization and simulation pipeline used for tree expansion. 
Compatible search statistics from the previous planning step are retained.
% TODO: Specify the root mutation operator and candidate budget in the implementation details; the present draft preserves the existing description.

\begin{algorithm}[t]
    \caption{Batched receding-horizon MCTS for sequence planning}
    \label{alg:mcts}
    \begin{algorithmic}[1]
        \Require State $\state_k$, previous search tree $\tree$, batch size $B$
        \Ensure Next state $\state_{k+1}$ and updated tree $\tree$
        \State Initialize or update $\tree$ at $\state_k$ and evaluate root candidates
        \While{the search budget remains}
            \State $\batch\gets\texttt{SELECTBATCH}(\tree,B)$
            \State $\mathcal{E}\gets\texttt{EVALUATEBATCH}(\batch,\initpose,\poseopt,\simul)$
            \State $\tree\gets\texttt{BACKUP}(\tree,\batch,\mathcal{E})$
        \EndWhile
        \State $\action^*\gets\argmax_{\action\in\mathcal{A}_{\mathrm{val}}(\state_k)} Q(\node_{\action})$
        \State Execute $\action^*$ and obtain $\state_{k+1}$
        \State Promote its child to the root and retain compatible descendants
        \State \Return $\state_{k+1},\tree$
    \end{algorithmic}
\end{algorithm}

\noindent \textbf{Batched Selection.}
The selection stage constructs a batch $\batch$ of search paths and
expansion requests. Progressive widening controls the number of children
available at each node as its completed visit count increases. When a node
admits a new child, its path is selected for expansion; otherwise, selection
continues through its existing children using the PUCT
score~\cite{silver2017mastering}.
% TODO: Specify the progressive-widening rule and how pending expansion
% requests count toward the child limit, consistently with the implementation.
For a child $\node$ reached from $\node_p$ by action $a$, let $Q(\node)$ be
the mean return associated with that action, $\visit(\node)$ its completed
evaluation count, and $\tmpvisit(\node)$ its temporary visit count within
the current batch. We select the child maximizing
\begin{equation}
    \puct(\node) =
    \frac{\visit(\node)Q(\node)}{\visit(\node)+\tmpvisit(\node)}
    + cP(\node_p,a)
    \frac{\sqrt{\visit(\node_p)+\tmpvisit(\node_p)}}
         {1+\visit(\node)+\tmpvisit(\node)},
    \label{eq:puct}
\end{equation}
where $c>0$ controls exploration and $P(\node_p,a)$ is an action prior
derived from geometric and equilibrium scores.
% TODO: Define the prior normalization and the initialization convention for
% an unvisited child, including the zero-denominator case in Eq. (\ref{eq:puct}).
After each path is selected, the temporary visit count is increased at every
node along that path. These reservations diversify subsequent selections
within the batch. All paths are selected before physical evaluation begins;
the reservations are removed during backup.

\noindent \textbf{Parallel Expansion and Evaluation.}
For each expansion request at state $\state$, $\initpose$ samples an object
and an initial pose $(i,\hat q)$, and $\poseopt$ refines the pose to obtain
the candidate action $a=(i,q)$. The simulation module $\simul$ then evaluates
the provisional placement and returns the settled successor state
$\state'=\transition(\state,a)$ together with its physical validity and
simulation statistics. These outputs, along with the optimization results,
determine the reward described above and whether the candidate is admitted
as a valid child. The batch evaluations are independent and run in parallel.
The evaluation stage returns the candidate outcomes to the tree update stage,
keeping the physical computations separate from tree selection and backup.

\noindent \textbf{Leaf Value Estimation and Backup.}
We estimate the remaining value at a nonterminal leaf without additional
rollouts, since each rollout step would require another placement
optimization and simulation. For a leaf state $\state_l$ at depth
$\depth_l$, measured from the current root, we use
\begin{equation*}
    \hat V(\state_l)=\bar r\sum_{j=0}^{h_l-1}\gamma^j,
    \qquad
    h_l=\min\!\left\{\depth_M-\depth_l,|\inventory(\state_l)|\right\},
\end{equation*}
where $\depth_M$ is the maximum search depth and $\bar r$ is the heuristic
mean reward estimated from evaluated placements. Terminal and failed leaves
receive no additional bootstrap value.
% TODO: Specify which evaluations contribute to the estimate of bar r, and
% define the failure return and whether failed candidates remain in the tree.

For a valid evaluated path $(\state_0,a_0,\ldots,a_{l-1},\state_l)$,
backup starts with $G_l=\hat V(\state_l)$, or zero for a terminal leaf,
and recursively computes
\begin{equation*}
    G_t=r(\state_t,a_t)+\gamma G_{t+1},
    \qquad t=l-1,\ldots,0.
\end{equation*}
The child corresponding to action $a_t$ receives $G_t$, so its value
includes both the immediate placement reward and the estimated subsequent
return. Each completed backup increments $\visit$ and updates $Q$ as an
empirical mean along the selected path. Temporary visit reservations are
then released; they do not themselves contribute samples to $Q$ or $\visit$.

\noindent \textbf{Action Selection and Tree Reuse.}
After the search budget is exhausted, the planner selects the root action
with the largest mean return among physically validated candidates:
\begin{equation*}
    a^*=\argmax_{a\in\mathcal{A}_{\mathrm{val}}(\state_k)}Q(\node_a),
\end{equation*}
where $\mathcal{A}_{\mathrm{val}}(\state_k)$ contains root actions with at
least one physically validated evaluation, and $\node_a$ is the child
corresponding to $a$. Candidates that fail pose optimization, simulation,
or feasibility checks are excluded from this set.
% TODO: Specify termination or recovery when the validated action set is empty.
Only $a^*$ is executed. Its child is promoted to the root, preserving
compatible descendants and their visit counts and return statistics.
Branches inconsistent with the observed settled state are discarded or
revalidated. The planner then repeats the search from the updated state
with a finite lookahead horizon.
